\documentclass[12pt]{ctexart}

\usepackage[a4paper,top=2.8cm,bottom=2.5cm,left=3cm,right=3cm]{geometry}
\usepackage{fontspec}
\usepackage{setspace}
\usepackage{enumitem}
\usepackage{booktabs}
\usepackage{longtable}
\usepackage{xurl}
\usepackage[colorlinks=true,linkcolor=black,urlcolor=blue,citecolor=black]{hyperref}
\usepackage[backend=biber,style=verbose-note,autocite=footnote,sorting=nyt]{biblatex}
\addbibresource{references.bib}

\setCJKmainfont[AutoFakeBold=3]{標楷體}
\setmainfont{Times New Roman}
\onehalfspacing
\setlength{\parindent}{2em}
\setlength{\parskip}{0.4em}
\setlist[itemize]{topsep=0.2em,itemsep=0.2em}
\setlist[enumerate]{topsep=0.2em,itemsep=0.2em}

\title{從副署爭議再論我國政府體制之爭\\
\large 行政院長究竟向誰負責？}
\author{王逸帆}
\date{\today}

\begin{document}
\maketitle

\section{前言：副署爭議的制度重量}

近期行政院長不副署財政收支劃分法修正案之爭議，使憲法第 37 條的副署制度重新進入政府體制討論。總統公布法律須否經行政院長副署，牽動文書程序，也牽動行政院長的責任歸屬：他究竟是在為總統政治團隊承擔責任，還是在責任政府架構下代表行政院向立法院負責。

本文從副署爭議出發，追問三個問題。第一，副署制度原本預設的是何種行政責任結構？第二，歷次修憲如何改變總統、行政院長與立法院之間的責任鏈？第三，當行政院長以不副署阻止法律公布時，這是責任政府的正常作用，還是總統制化後的第二道行政否決權？

\section{問題意識：政府體制分類與制度效果}

我國中央政府體制之爭，經常以總統制、半總統制或雙首長制等分類展開。這些分類有其說明功能，卻容易停在名稱選擇。副署爭議使制度效果變得無法迴避：行政院長能否以不副署方式，使立法院已三讀且覆議後仍維持的法律無法公布？若答案為肯定，副署便具有接近行政否決的效果；若答案為否定，憲法第 37 條的副署要求又須被重新說明。

本文以副署制度為入口，檢驗憲法本文與增修條文之間的責任配置。憲法第 37 條仍保留副署制度，增修條文第 3 條仍稱行政院對立法院負責；同時，行政院長由總統任命，總統發布行政院長任免命令也無須行政院長副署。這幾個規範放在一起，產生一個必須處理的問題：行政院長在法律公布程序中的副署權，究竟是在承擔政治責任，還是在形成另一道行政權否決。

\section{研究方法：由指定書目建立問題群}

本文的方法起點，是憲法解釋教材所提出的問題群。文義、歷史、體系與目的解釋仍是後續分析材料，但不作為固定框架。副署爭議牽涉條文語句、修憲時間、機關結構、政治責任、憲法審查角色與制度後果；若從單一方法進入，容易過早把問題縮成第 37 條幾個字的法律效果。本文因此先建立書目矩陣，再依矩陣整理事實與材料。

\begin{longtable}{p{3.1cm}p{4.2cm}p{6.7cm}}
\toprule
書目群 & 方法問題 & 對副署研究的要求 \\
\midrule
BF & 憲法解釋是否能避開民主、機關角色與制度風險判斷 & 檢查第 37 條文字、責任政府概念、總統直選、行政院長任命與法律公布效果之間的關係。\autocite[26--30]{barber_fleming_2007}\\
GAF / Bobbitt & 歷史、文義、結構、先例、倫理與審慎等論證模式如何並用 & 對每一種副署定位，分別列出其文字負擔、修憲史負擔、結構負擔與制度後果負擔。\autocite[4]{bobbitt_gaf_5ed}\\
LA / Perry & 憲法文本、憲法規範與具體化判斷如何區分 & 先辨認「行政院對立法院負責」在現制下是否仍是可操作規範，再討論由誰具體化。\autocite[99--151]{perry_1998}\\
LA / Rubenfeld & 憲法承諾如何跨時間延續 & 把 1947 年本文、1994 年總統直選、1997 年行政院長任命改制與近年副署實踐放在同一時間軸檢查。\autocite[194--235]{rubenfeld_1998}\\
DB & 法治、比例與國家行為理由如何連結 & 比例公式不直接移植到機關爭議；其可用之處，在於檢查不副署理由的強度、程序與停滯效果。\autocite[159--188]{beatty_2004}\\
一般法學方法 & 文義邊界、體系衝突、漏洞填補與法續造界線如何處理 & 若將副署理解成阻斷法律公布的權限，須說明此一效果是否仍在第 37 條可承載範圍內。\autocites{yang_jurisprudence_2017}{lee_methodology_2018}{larenz_methodology_2008}\\
憲法解釋原則 & 解釋方法與解釋原則如何區分 & 以節制解釋權、忠於憲法結構、重視比較法與制度後果作為檢驗尺度，避免由機關自行決定副署效果。\autocite[559--627]{chen_constitutional_interpretation_2019}\\
政府體制與權力分立 & 政府體制案件如何處理責任鏈與僵局出口 & 將副署、覆議、不信任、解散、行政院長任命與憲法審查作為同一組制度材料。\autocites{tang_constitutional_structure_2014}{tang_judicial_review_balance_2014}{tang_tiered_proportionality_2009}\\
\bottomrule
\end{longtable}

這個矩陣使本文保留三條研究要求。第一，副署的文字效果要和修憲史一併檢查。憲法第 37 條仍以「須經副署」作為一般規則；增修條文第 2 條列出免副署事項；增修條文第 3 條又保留行政院對立法院負責、覆議、不信任與解散等制度。這些規範共同構成研究對象。

第二，先例與本土文獻只能在可移植範圍內使用。林春元對政黨政治下權力分立的分析，提供觀察一致政府、分立政府、少數政府與政黨極化的欄位。\autocite[325--370]{lin_chunyuan_2013} 賈文宇對「前解釋」的研究，提醒本文引用釋字第 387、419、499 號時，須區分作成時間、規範問題、理由結構與現制差異。\autocite[613--641]{jia_wenyu_2021} 蔡季廷關於美國總統簽署聲明書的研究，則提供行政部門對國會法律提出憲法異議時，聲明、否決與法律生效之間如何區分的比較材料。\autocite[85--146]{tsai_jiting_2013}

第三，比較法不作制度名稱的移植。臺大開放式課程「大法官講堂：中華民國憲法及政府」在總統與行政院講綱中，已把副署、總統任命、覆議、不信任與解散放在政府體制脈絡中處理。\autocites{ntu_ocw_separation_lecture}{ntu_ocw_executive_yuan_lecture} 德國材料用來檢查副署與議會責任政府如何連動；法國材料用來觀察雙元行政、免副署清單與同居政治；美國材料用來對照明文否決權、覆決門檻與行政異議聲明。這些比較材料只提供問題清單，不預先決定我國第 37 條的答案。

本文現階段只提出研究設計。後續須完成三項工作：其一，逐條整理 1994、1997 年修憲紀錄與官方文本；其二，建立 2025、2026 年不副署案例表，欄位包括時間、標的、總統、行政院長、立法院席次、朝野關係、覆議狀態、不副署理由與後續程序；其三，把釋字材料、憲法法庭判決與比較法制度放回上述矩陣中檢驗。副署是責任標示、程序生效要件、合憲性把關機制，或接近行政否決的工具，須待材料對讀完成後再判斷。

\section{副署制度的原始功能}

\subsection{憲法第 37 條的責任政府前提}

憲法第 37 條規定，總統依法公布法律、發布命令，須經行政院院長之副署，或行政院院長及有關部會首長之副署。\autocite{roc_constitution_moj} 這個制度若放在憲法本文的整體結構中理解，已超出文書簽名程序。副署的原始功能，是將總統的形式行為連回可負政治責任的行政院長或部會首長，使總統不直接與聞日常行政，並使行政責任能被立法院追問。

這一點必須和憲法第 55 條、第 57 條合併理解。憲法第 55 條原本規定行政院長由總統提名、經立法院同意任命；第 57 條則規定行政院對立法院負責。行政院長在此結構中同時連結總統提名與立法院同意。副署制度也在這個安排下取得意義：總統的公布法律與發布命令，經由行政院長副署而接上國會責任政治。

\subsection{副署與實質否決權的區分}

若副署被理解為責任標示，行政院長的角色重在承擔政治責任，不宜逕行推導出任意阻止公布法律的權限。相反地，若將副署理解為行政院長可任意拒絕的生效要件，副署就可能從責任政府工具轉變成行政權阻斷法律公布的制度武器。後續研究須進一步檢查：副署在包含審查與政治負責功能時，是否仍須接上請辭、倒閣、解散或憲法審查等責任政治出口。

\section{修憲如何改變副署與責任鏈}

\subsection{七次修憲中與本題相關者}

中華民國憲法自 1991 年起共有七次增修。與副署及政府體制最直接相關者，狹義是 1994 年第三次增修與 1997 年第四次增修；廣義則還應納入 1991 年第一次增修與 1992 年第二次增修，因其分別處理總統國安實權與總統直選民主正當性。

\begin{longtable}{p{2.2cm}p{3.2cm}p{8cm}}
\toprule
時間 & 修憲 & 本題意義 \\
\midrule
1991 & 第一次增修 & 終止動員戡亂後，仍保留總統決定國家安全大政方針的制度位置，總統實權化沒有完全退回憲法本文的形式元首想像。\\
1992 & 第二次增修 & 明定總統、副總統由自由地區人民選舉，總統直接民主正當性的制度基礎開始形成。\\
1994 & 第三次增修 & 具體化總統直選，並限縮副署：總統發布部分經國民大會或立法院同意任命人員之任免命令，無須行政院長副署。\\
1997 & 第四次增修 & 行政院長由總統任命，停止適用憲法第 55 條；總統發布行政院長任免命令及解散立法院命令，無須行政院長副署；同時保留行政院對立法院負責。\\
2000、2005 & 第六、七次增修 & 調整國民大會與修憲複決制度，副署與行政院長任免結構基本延續 1997 年設計。\\
\bottomrule
\end{longtable}

\subsection{1994 年第三次增修：副署限縮的開端}

1994 年第三次增修具體化總統直選制度，並同時規定總統發布依憲法經國民大會或立法院同意任命人員之任免命令，無須行政院院長副署。\autocite{president_third_amendment} 這一變化的意義，在於總統的人事權開始從行政院長副署中部分脫離。

此時的副署限縮尚未徹底改造行政院長的權力來源，因為行政院長本身仍有立法院同意基礎。然而，總統直選與副署限縮已經共同預示後來的方向：總統不再只是形式國家元首，而逐步取得直接民主正當性與更大的實質人事空間。

\subsection{1997 年第四次增修：責任鏈的斷裂點}

1997 年第四次增修是本文最重要的制度節點。增修條文第 3 條規定行政院院長由總統任命，停止適用憲法第 55 條。這使行政院長不再須經立法院同意，其權力來源明顯轉向總統。同時，增修條文第 2 條又規定總統發布行政院院長任免命令及解散立法院命令，無須行政院院長副署。\autocite{president_fourth_amendment} 行政院長既由總統任命，其任免命令又不受行政院長副署制衡，副署作為行政院長制衡總統的功能因此大幅弱化。

但 1997 年修憲並未刪除「行政院對立法院負責」的語彙，施政報告、質詢、覆議、不信任投票與解散立法院也被重新編排於增修條文第 3 條。矛盾由此形成：行政院長的產生與存續轉向總統，但行政院對立法院負責的文字仍被保留。

\section{副署爭議的現實案例群}

\subsection{1991 年郝柏村副署爭議}

政治史上常被提及的第一個副署爭議，是 1991 年行政院長郝柏村與李登輝總統之間的人事命令副署爭議。此一事件常被稱為拒絕副署前例，但其事實型態仍須區分：若當時並無正式人事命令送請副署，而僅屬政治上表達反對，則不宜逕視為嚴格意義的不副署。即使將其視為副署爭議，它的性質也主要是總統與行政院長之間的人事權衝突，尚非立法院已維持法律案後的法律公布問題。

\subsection{2025 年財劃法不副署爭議與反年改對照案}

2025 年財政收支劃分法修正案引發的不副署爭議，憲政意義更為尖銳。爭議已經進入法律公布階段：立法院通過法律案，行政院行使覆議程序；若覆議失敗後行政院長仍拒絕副署，法律將無法公布。行政院相關說明主張不副署是為維護憲政秩序與權力分立；媒體亦將之描述為法律公布層次的不副署憲政首例。\autocites{ey_finance_countersign_2025}{initium_2025_finance_countersign}

財劃法案的不副署理由，可以先分成兩層觀察。總統府批示部分，重點在中央財政削弱、中央地方分配不均、財政永續、民生政策排擠，以及國安與災變能力受阻。行政院長加註部分，則以憲法第 37 條副署要求為形式基礎，推導行政院長得不予副署，並主張修法侵害行政院預算編列權、悖離民主原則與權力分立，且違憲情形重大明顯。\autocites{president_finance_2025}{ey_finance_unconstitutional_pdf_2025} 這些理由後續須分別檢查：財政衝擊是否足以支持不副署，預算編列權與立法院法律形成權如何分界，以及行政院長能否自行判斷違憲明顯性。另須注意，「不副署權」一語未出現在憲法第 37 條本文；此處應理解為行政院從副署制度中提出的權限主張。

行政院後續院會資料又提供另一層定位。第 3983 次院會將行政院長依第 37 條所採取的措施，放在憲政體制維護、憲政秩序與憲政機關互動中說明；第 3984 次院會則稱行政院長代表政府作成不副署決定，目的在維護憲政秩序運作，並促使朝野政黨重新協商。\autocite{ey_meeting_153_dataset} 這表示財劃法案中的不副署理由，已包含法案違憲判斷、行政權限防衛與促成政治協商三個層次。

同年反年改相關修法則提供一個重要對照。行政院面對其認為有憲法疑義的法案時，也曾選擇副署後聲請憲法審查。\autocite{ey_pension_review_2025} 這使本文後續研究不能只問行政院長是否有拒絕副署權，還要比較不副署與副署後憲法審查兩種路徑，各自如何影響法律公布、機關責任與憲政僵局。

\subsection{2026 年續發不副署案例}

2026 年 2 月，國軍老舊眷村改建條例修法又發生不副署爭議。依中央社報導，行政院與行政院長說明不副署理由時，指向特定個案修法、平等原則、司法判決與行政預算權等疑義。\autocites{cna_village_2026_1}{cna_village_2026_2} 同年 3 月，行政院又對黨產條例、衛星廣播電視法、立法院組織法三項修法決定不予副署。\autocite{cna_three_bills_2026}

2026 年續發案例的理由型態，比財劃法更分散。眷改條例涉及特定個案修法、平等原則、既有司法判決、個案立法禁止原則、基金負擔，以及行政院施政計畫與預算形成權。三法不副署則分別指向衛星廣播電視法的頻位資源與差別待遇、黨產條例對既有合憲判斷與特定個案的處理、立法院組織法的助理費制度、公開透明、財政負擔與法治國原則。這些理由的共同特徵，是行政院提出的理由已從單一財政或程序問題，擴及平等、權力分立、司法判決效力、財政負擔與民主程序等多組理由。

這些續發案例使副署問題超出財劃法單一事件。待研究的重點包括：不副署是否正在成為行政院處理高爭議國會法案的常規工具；行政院在不副署與副署後聲請憲法審查之間作選擇時，應提出何種程度的憲法理由；若不副署反覆發生，立法院、總統府與憲法法庭是否仍有足夠程序解除僵局。

\section{後續研究問題}

本文此階段不提出副署權的最後定位。研究仍須完成課本與本土文獻的對讀，再決定正文的論證方向。暫定處理下列問題。

\subsection{憲法解釋的對象：條文、制度，或憲政實踐}

第 37 條副署、第 55 條行政院長任命、第 57 條行政院對立法院負責，以及增修條文第 2、3 條，不能被拆成孤立條文處理。待研究的問題是：本題的解釋對象究竟是「副署」這一個字面要件，還是總統、行政院長與立法院之間的責任配置。BF 所提出的「憲法是什麼」與「誰有權解釋憲法」兩個問題，應先放在這裡處理。

\subsection{時間問題：制憲本文與增修條文如何共同拘束}

憲法本文保留責任政府設計，增修條文則逐步強化總統直選與行政院長任命權。待研究的問題是：1997 年修憲究竟改變了責任政府的哪些部分，又保留了哪些部分。GAF、LA 與 Rubenfeld 關於憲法承諾、歷史與當代理由的討論，應用來檢查這個時間斷裂，避免直接把修憲後政治效果當成答案。

\subsection{政府體制案件中的解釋者角色}

副署爭議若進入憲法審查，解釋者面對的是憲法機關之間的權限與責任分配，審查方式未必能沿用一般基本權案件。待研究的問題是：法院或憲法法庭在此類案件中應只定分止爭，還是可以補足政治部門未能形成的責任出口。這一部分須閱讀權力分立、政府體制與憲法審查角色相關文獻。

\subsection{制度後果：副署、不副署與憲政僵局}

不副署的效果不能先被假定。待研究的問題包括：不副署是否使法律公布程序暫停；行政院長是否負說理義務；覆議失敗後能否再以不副署阻止公布；若立法院、行政院、總統府與憲法法庭對此意見不同，僵局如何解除。Beatty 關於法治與比例的討論、GAF 中程序論與民主理論的材料，可用來檢查不同方案的制度成本。

\subsection{比較對象：德國、法國與總統制對照}

比較法研究暫以德國、法國與美國為主要對照。德國基本法第 58 條將聯邦總統命令與指示原則上繫於聯邦總理或主管部長副署；同時，第 63、67、68 條又把總理選任、建設性不信任投票與信任案解散放進議會責任政府架構。\autocite{germany_basic_law} 德國材料可用來檢查一件事：副署若要承擔責任政治功能，是否需要同時具備清楚的政府首長產生方式、去職方式與僵局出口。

法國第五共和憲法提供另一組問題。第 8 條規定總統任命總理，第 12 條規定總統得解散國民議會，第 19 條則明列若干總統行為不須副署；第 20、49、50 條又規定政府決定並執行國家政策，且向國會負責。\autocite{elysee_constitution} 法國材料應用來觀察雙元行政、反副署事項與同居政治的制度配套，尤其是總統多數與國會多數不一致時，總理如何取得政策與責任位置。

美國則作為總統制的程序對照。美國憲法第一條第七款把總統同意、退回異議、兩院三分之二重新通過，以及十日期限明文規定。\autocite{loc_us_constitution_article1} 這項材料可協助追問：若不副署在我國被理解為接近行政否決權的效果，憲法解釋是否需要同時處理理由提出、覆決程序與法律公布僵局。

三個比較對象的功能不同。德國重在副署與議會責任政府，法國重在雙元行政與反副署清單，美國重在明文否決權與國會覆決。本文後續研究將以這些制度配套作為問題清單，再回到我國憲法第 37 條與增修條文第 2、3 條的結構中判斷。

\section{待檢驗命題：副署爭議中的四組問題}

\subsection{第 37 條文字與例外清單}

第 37 條「須經副署」具有強制語氣。增修條文第 2 條又列出若干無須副署事項，包括行政院長任免命令、依憲法經立法院同意任命人員之任免命令，以及解散立法院命令。由此可形成第一個待檢驗命題：法律公布仍在第 37 條副署要求之內，但副署要求的法律效果尚待說明。

這個命題需要回答三個問題。其一，副署欠缺時，總統是否完全不能公布法律。其二，行政院長拒絕副署時，是否必須同時提出可受檢驗的憲法理由。其三，若拒絕副署造成公布程序停滯，憲法秩序中是否存在時間限制或救濟出口。

\subsection{制憲本文、增修條文與憲法承諾}

制憲本文中，第 55 條、第 57 條與第 37 條共同形成責任政府語境：行政院長經立法院同意任命，行政院對立法院負責，總統行為經行政院長副署而連回政治責任。1994 年修憲開始限縮副署，1997 年修憲又使行政院長改由總統任命。第二個待檢驗命題是：行政院長權力來源已轉向總統，但行政院對立法院負責的憲法語彙仍須被解釋。

這裡的研究重點在於檢查 1947 年本文與 1997 年修憲如何共同限制今日副署制度。Rubenfeld 的憲法承諾論可協助處理這個問題：若「行政院對立法院負責」仍是一項跨時間承諾，副署制度在 1997 年後應如何保留責任政治意義；若該承諾已大幅轉化，轉化的界線與理由也需要說明。

\subsection{釋字材料與先例使用}

釋字第 387 號可作為憲法本文時代責任政府的官方材料，因其明示行政院長對立法院負政治責任，並將總統、行政院與立法院放入制度性設計中。釋字第 419 號則顯示政府體制問題不能只處理單一條文，還要檢查職位相容性、權力分立與制衡功能。釋字第 499 號可作為權力分立、制衡與政治責任仍屬憲法整體原則的背景材料。

賈文宇對「前解釋」的分析提醒，本文引用這些釋字時，必須區分可直接使用的命題與只能作為背景的材料。釋字第 387、419 號均早於 1997 年第四次修憲，不能把其結論直接移到現行增修條文；釋字第 499 號處理的是第五次修憲，也不能被寫成已經處理 1997 年政府體制設計。本文後續須逐一標示各釋字的時間、問題、理由與可用範圍。

\subsection{不副署理由的類型化}

2025 年財劃法案中，官方理由包括中央財政削弱、中央地方分配不均、預算編列權、權力分立、民主原則、財政永續與憲法忠誠。2026 年眷改條例與三法不副署案例，理由又擴及特定個案修法、平等原則、司法判決效力、頻位資源、黨產處理、助理費制度與財政負擔。這些理由不能合併成「維護憲政秩序」一句處理。

本文後續將把理由分為五類：副署權限理由、行政權與預算權理由、民主程序理由、平等與個案立法理由、司法判決與釋憲結果理由。每一類理由都要回答同一組問題：該理由是否具有憲法位階，是否足以阻卻法律公布，是否有較低衝突的替代程序，行政院長提出該理由後需承擔何種政治責任。

\section{比較法材料的研究功能}

\subsection{德國：副署與議會責任政府}

德國基本法第 58 條將聯邦總統命令與指示原則上繫於聯邦總理或主管部長副署；第 63、67、68 條則安排總理選任、建設性不信任投票與信任案解散。\autocite{germany_basic_law} 德國材料的用處，在於觀察副署如何嵌入清楚的議會責任政府。若副署要承擔責任歸屬功能，政府首長的產生、去職與僵局出口需要同時被看見。

\subsection{法國：雙元行政與反副署清單}

法國第五共和憲法第 8 條規定總統任命總理，第 12 條規定總統得解散國民議會，第 19 條明列若干總統行為免副署；第 20、49、50 條則規定政府決定並執行國家政策，且向國會負責。\autocite{elysee_constitution} 法國材料的用處，在於觀察雙元行政如何透過免副署清單、總理責任與同居政治取得運作空間。臺灣若借用半總統制語彙，必須同時比較這些配套，而不能只比對總統直選與總理任命。

\subsection{美國：行政異議與明文否決程序}

美國憲法第一條第七款規定總統同意、退回異議、兩院三分之二重新通過，以及十日期限。\autocite{loc_us_constitution_article1} 蔡季廷對美國總統法案簽署聲明書的研究則顯示，行政部門對法律提出憲法異議時，其效果可能只是輔助解釋或政治表態，未必阻止法律生效。這個比較對象可用來檢查臺灣不副署若具有阻斷公布效果，是否已接近明文否決權；若接近，程序、期限與覆決機制又如何補足。

\section{下一階段研究安排}

本文尚不提出副署權的最終定位。下一階段將集中處理四項工作。第一，查核 1994、1997 年修憲紀錄中關於副署、行政院長任命、覆議與不信任投票的討論。第二，補齊 2025、2026 年不副署案例的正式文書、立法院程序與後續救濟。第三，逐一整理釋字第 387、419、499 號及近年憲法法庭判決在政府體制案件中的方法。第四，將德國、法國、美國制度與臺灣現制放在同一張比較表中，確認哪些比較材料屬於背景，哪些足以幫助說明第 37 條副署的法律效果。

\printbibliography

\end{document}
